\section{Known Limitations}
\label{sec:known-limitations}

This section states each limitation explicitly rather than letting any of them evaporate under
space pressure. None is new information relative to Section~\ref{sec:results}; the point of
collecting them here is that a reader assessing how far to trust this paper's claims should be
able to find every scope boundary in one place.

\paragraph{No live field deployment; this is an offline, bench-level evaluation.} This is the most
foundational scope boundary in the paper, stated first rather than buried among smaller caveats. No
experiment reported in this paper was conducted with real physical sensors or inside an operating
recirculating aquaculture facility -- no physical water-quality sensor hardware exists for this
system at all. Every measurement -- the vision-model validation, the RAG retrieval tests, the
AQUA-1B/rule-engine behavioral trace, the agentic tool-routing benchmark, the PSI drift-binning
sweep, and the edge export benchmark -- was run against synthetically generated sensor data, static
labeled images, offline analysis scripts, and scripted or synthetic scenarios executed against the
codebase, not against live sensor telemetry from an operating facility. This is a deliberate choice
of study design, not an incidental gap: absent physical sensor hardware, the goal of this evaluation
is to demonstrate the architectural feasibility of the proposed multi-component pipeline and to test
whether its space and memory footprint is viable on realistic edge hardware, rather than to validate
decision-making accuracy in a live setting; Section~\ref{sec:introduction} and
Section~\ref{sec:experimental-setup} state this framing directly. Only the system's interface has
ever been demonstrated, as a mockup/demo; the system itself has not been field-tested. Separately,
and just as plainly: S.U.R.E. was built for the TEKNOFEST Agricultural Technologies competition, but
this team did not enter that competition and did not pass its preliminary evaluation round. Every
claim elsewhere in this paper about what the system does or how it performs should be read against
this scope, not against an assumption of field deployment.

\paragraph{Sample size of the central result.} The four-bucket dual-layer classification
(Section~\ref{sec:results-dual-layer}) is measured on $n=8$ -- the evaluation harness's entire
fixed scenario population, not a sample drawn from a larger population. No confidence interval is
implied by this count and none should be inferred from it. The finding is a census of a small,
fixed scenario set, not a statistically powered estimate of a rate that would hold over a larger or
different set of scenarios.

\paragraph{Single-device edge benchmarking.} Every latency and export-format number in
Section~\ref{sec:results-vision} was measured on a single Apple M4 Pro. The CoreML/ANE figure was
re-measured across three independent process invocations to rule out single-session variance, but
multi-session thermal-throttling effects beyond those three sessions, and behavior on any other
device class, were not characterized and should not be assumed to transfer.

\paragraph{twin\_bridge remains unresolved.} The digital-twin comparison module is unit-tested
only. It has never been exercised against a live Godot or CODESYS session, and no such session
exists anywhere in this project's history. This paper never describes it as field-validated, and a
reader should not infer field validation from the passing unit-test count.

\paragraph{The adapter-provenance caveat is scoped narrowly.} It applies to results derived from
the dual-layer decision system (Section~\ref{sec:results-dual-layer}), which used the configured
\texttt{sure-aqua-adapter}. It does not apply to the agentic tool-routing results
(Section~\ref{sec:results-agentic}), which ran the base model with no adapter loaded at all.
Applying this caveat to the agentic-routing result would be a new factual error, not added
precision, and this paper does not make that error anywhere in its text.

\paragraph{The vision dataset is real but small.} Unlike the sensor values discussed above, the
510-image detector-training dataset (Section~\ref{sec:vision-pipeline}) is composed of genuine
photographs captured from the project's own physical camera setup and hand-labeled by the team, not
synthetic imagery. It is, however, intentionally small at this stage, and the authors plan to expand
it as future work; no claim in this paper should be read as asserting that 510 images is a
sufficient long-run dataset size, only that it is what the reported detector was actually trained
and validated on.

\paragraph{The validation set has no sparse-fish frames.} Every one of the 98 validation frames
used throughout Section~\ref{sec:results-vision} contains at least three ground-truth fish. The
reassuring 0/98 full-frame-miss result is conditional on that property and says nothing about the
$k=1$/$k=2$ regime, which this dataset does not contain at all. The 22--28\% miss-risk figure
sometimes cited alongside it is an extrapolation from an independence approximation, not a
measurement, and should be read as one.

\paragraph{A live, unfixed data-integrity issue.} The system's RAG knowledge base
(\url{llm-service/knowledge/06-davranis-ve-refah-gostergeleri.md}) still states the detector's
recall as approximately 0.695 as a live fact, a figure superseded elsewhere in this system's
history 43 minutes after it was introduced, but never corrected in this specific file. This file is
ingested into the system's vector store and is retrievable by the LLM running against it. It was
not fixed during this audit, consistent with a read-only guardrail on the running system; it is
disclosed here as a genuine, currently unresolved issue in the codebase, not a historical footnote.

\paragraph{Minor, non-conclusion-changing inconsistencies.} A thousandths-place rounding
difference exists between the commonly reported precision figure (0.858) and a freshly reproduced
raw \texttt{val()} output (0.859); neither number changes any claim in this paper. AQUA-7B's
published agent-benchmark format-compliance and mean-step-count figures did not reproduce exactly
on a fresh, code-unmodified re-run, with no configuration difference found to explain the drift,
though the selection-accuracy and constant-answer conclusions reproduced exactly. The PSI
drift-binning sweep's own $\delta=0.04$ entry differs from a separately drawn significant-drift
window at the same nominal severity, an artifact of unseeded per-call resampling; no headline
number in this paper depends on either individual value.
